\begin{figure}[!htbp]
\centering
\begin{tikzpicture}
  \node[grey, minimum width=30mm] (S) at (0,0) {\textbf{Sampling}};
  \node[slate, minimum width=56mm] (P) at (-3.6,-1.4) {\textbf{Probability}\\\footnotesize every unit has a known chance};
  \node[sand, minimum width=56mm] (N) at (3.6,-1.4) {\textbf{Non-probability}\\\footnotesize chance of selection unknown};
  \draw[arr] (S) -- (P); \draw[arr] (S) -- (N);
  \node[plain, text width=54mm, align=left, font=\footnotesize, anchor=north] at (-3.6,-2.2)
    {\textbf{simple random} — lottery, random tables\\\textbf{systematic} — every $k$-th unit\\\textbf{stratified} — random within homogeneous strata\\\textbf{cluster} — whole natural groups\\\textbf{multi-stage} — clusters, then units};
  \node[plain, text width=54mm, align=left, font=\footnotesize, anchor=north] at (3.6,-2.2)
    {\textbf{convenience} — whoever is available\\\textbf{purposive} — chosen by judgement\\\textbf{quota} — fill fixed numbers per group\\\textbf{snowball} — referral chains to hidden groups\\\phantom{x}};
\end{tikzpicture}
\caption{Sampling methods. Only probability sampling allows the sampling error to be estimated and the
results to be generalised statistically.}
\end{figure}
